\section{Verification purpose}
This case verifies live-fuel moisture damping and calculated live extinction while fuel loads remain fixed.  The static shrub models SH2, SH6, and SH8 are intentionally used so changing live moisture cannot transfer herbaceous load into a dead class.  It can expose live/dead index errors, percent conversion errors, incorrect live-extinction calculation, or unintended dynamic curing.  The question is whether all static-fuel variants match independent $I_R$ and ROS predictions over \LiveMoistureSet{}.

\section{Mathematical formulation and numerical implementation}
The live term uses the same cubic damping function $\eta_M(r_M)$ as dead fuel, but $r_M=M_L/M_{x,L}$ and $M_{x,L}$ is calculated from dead/live weighting and dead moisture rather than read directly.  Live moisture also appears in $Q_{ig}=250+1116M_f$.  The formulation follows Andrews (2018), RMRS-GTR-371.  For a static fuel, $W_{1h}$, $W_{LH}$, and $W_{LW}$ are invariant throughout the sweep.

\section{Derivation of expected behavior}
The oracle first calculates the live extinction moisture for each frozen fuel row and fixed dead moisture.  It then evaluates live damping, reaction intensity, heat sink, and $R_0$ at each moisture.  Since no loads move, changes arise only from moisture physics.  This controlled result is the baseline needed to interpret the additional load-transfer mechanism in CASE39.

\section{Verification metrics and acceptance criteria}
Maximum relative errors are 0.02 for direct ROS and 0.01 for reaction intensity.  Variants predicted at or above \SI{0.10}{m.min^{-1}} also require TOA-derived ROS error at most 0.05 and minimum $R^2=0.995$; slower variants use the direct field because the practical run-duration cap cannot provide a multi-cell TOA regression.  Pointwise analytical comparison is selected because a merely monotonic curve would not detect an incorrect live-extinction threshold.  All 24 variants are required and the preprocessor rejects a fuel row marked dynamic.

\section{Simulation configuration and predicted outcome}

The finite-strip PHI raster is a signed-distance field clipped to ELMFIRE's supported $[-1,1]$ interval, preserving both an ignited interior and a resolved front gradient. For every positive oracle ROS, \texttt{SIMULATION\_DT} and \texttt{SIMULATION\_DTMAX} are selected so the expected front advances 5\% of one cell per step (CFL 0.05, below the configured 0.2 limit). This keeps the PHI change above ELMFIRE's fixed stall scale even if the implementation produces only a small fraction of the reference speed. A zero-ROS variant retains a 5 s step and is evaluated from direct fields without requiring TOA propagation. The analytically selected stop time is rounded upward to a whole second because the current ELMFIRE time loop stores the duration as an integer; this preserves the intended duration and guarantees the terminal output dump. The evaluator always requires uniquely named rasters tied to the one final dump record. A TOA-required variant must reach the \texttt{SIMULATION\_TSTOP} in its generated namelist; a direct-only variant may terminate early because lack of propagation is part of the state being tested.
Live-herb and live-woody moisture are both set to 30, 50, 70, 90, 120, 150, 180, or 200\%; dead classes remain at 8\%.  The homogeneous domain uses a finite ignition strip whose east face is planar in the interior measurement corridor and whose ends remain inside the lateral boundaries.  Wind and slope are zero and all coupled processes are disabled.  Fuel loads remain invariant; reaction intensity and ROS should generally decline as live moisture rises, with behavior governed by each fuel's calculated live extinction.

\section{Actual simulation results}
The generated record contains \CompletedVariantCount{} of \RequiredVariantCount{} required variants and reports \textbf{\VerificationStatus}.
\ifnum\CompletedVariantCount>0
\IfFileExists{../figures/sweep_response.pdf}{\begin{figure}[H]\centering\includegraphics[width=0.95\textwidth]{sweep_response.pdf}\caption{Hand calculations (open diamonds, solid lines) and ELMFIRE predictions (filled circles, dashed lines) for the static live-moisture response. Each fuel-model pair uses the same color.}\end{figure}}{}
\fi

\section{Calculated metrics and verification decision}
\begin{center}\begin{tabular}{@{}llll@{}}\toprule Metric & Limit & Calculated & Status\\\midrule\MetricRows\bottomrule\end{tabular}\end{center}
The authoritative decision is \textbf{\VerificationStatus}.

\section{Technical interpretation}
A shifted transition implicates calculated live extinction or moisture units.  Unexpected sharp changes suggest that a supposedly static row was treated dynamically.  Agreement in intensity but not ROS narrows the defect to heat sink or spread conversion.

\section{Scope and limitations}
This case covers three static shrub fuels at one dead-moisture state.  It does not verify dynamic curing, spatial moisture mosaics, temporal interpolation, or physical observations.  Dynamic herbaceous transfer is isolated in CASE39.
